\section*{Chapter \ref{ch:b1:faults-folds-fabrics} --- Faults, Folds and Fabrics}

\begin{solution}{exo:b1:faults-folds-fabrics:1}
(a) A \omterm{def:b1:faults-folds-fabrics:fault-types}{normal fault}. (b) A \omterm{def:b1:faults-folds-fabrics:fault-types}{reverse fault}, nearly pure dip-slip. (c) A \omterm{def:b1:faults-folds-fabrics:fault-types}{sinistral}
\omterm{def:b1:faults-folds-fabrics:fault-types}{strike-slip fault}. (d) A \omterm{def:b1:faults-folds-fabrics:fault-types}{dextral} \omterm{def:b1:faults-folds-fabrics:fault-types}{strike-slip fault}. (e) An \omterm{def:b1:faults-folds-fabrics:fault-types}{oblique-slip fault}:
the \omterm{def:b1:faults-folds-fabrics:fault-elements}{hanging wall} moved down and along the \omterm{def:b1:geological-maps:attitude}{strike}, to the left; its strike-slip
component, $\cos\ang{40} = 0.77$ of the \omterm{def:b1:faults-folds-fabrics:fault-elements}{net slip}, is larger than its dip-slip
component, $0.64$, so it is called a \omterm{def:b1:faults-folds-fabrics:fault-types}{sinistral} fault with a normal component.
(f) With a \omterm{def:b1:geological-maps:attitude}{dip} of \ang{25}, less than \ang{45}, the \omterm{def:b1:faults-folds-fabrics:fault-types}{reverse fault} is a \omterm{def:b1:faults-folds-fabrics:fault-types}{thrust
fault}.
\end{solution}

\begin{solution}{exo:b1:faults-folds-fabrics:2}
With $|\sin r| = 1$, the \omterm{def:b1:cross-sections:throw}{throw} is $300\sin\ang{55} = \qty{245.7}{m}$ and the
\omterm{def:b1:cross-sections:throw}{heave} $300\cos\ang{55} = \qty{172.1}{m}$.
\end{solution}

\begin{solution}{exo:b1:faults-folds-fabrics:3}
For an \omterm{def:b1:faults-folds-fabrics:fold-classes}{upright fold} whose \omterm{def:b1:faults-folds-fabrics:fold-elements}{limbs} \omterm{def:b1:geological-maps:attitude}{dip} $\delta_1$ and $\delta_2$ in opposite
directions, $\iota = \ang{180} - \delta_1 - \delta_2$. (a) \ang{110}: open.
(b) \ang{35}: tight. (c) \ang{7}: isoclinal.
\end{solution}

\begin{solution}{exo:b1:faults-folds-fabrics:4}
\omterm{def:b1:faults-folds-fabrics:fault-types}{Normal faults} \omterm{def:b1:geological-maps:attitude}{dip} $\ang{45} + \ang{35}/2 = \ang{62.5}$, thrusts
$\ang{45} - \ang{17.5} = \ang{27.5}$. Conjugate \omterm{def:b1:faults-folds-fabrics:fault-types}{strike-slip faults} each make
\ang{27.5} with $\sigma_1$, on either side of it, so they meet at \ang{55}.
\end{solution}

\begin{solution}{exo:b1:faults-folds-fabrics:5}
(a) The strike-slip component is $400\cos\ang{30} = \qty{346.4}{m}$ along the
\omterm{def:b1:geological-maps:attitude}{strike}, the dip-slip component $400\sin\ang{-30} = \qty{-200}{m}$, that is
\qty{200}{m} down the \omterm{def:b1:geological-maps:attitude}{dip}; the \omterm{def:b1:cross-sections:throw}{throw} is $200\sin\ang{50} = \qty{153.2}{m}$ and
the \omterm{def:b1:cross-sections:throw}{heave} $200\cos\ang{50} = \qty{128.6}{m}$. (b) The \omterm{def:b1:geological-maps:attitude}{strike} is $000$, so
$\mathbf s = (0, 1, 0)$; $\mathbf u = (-\cos\ang{50}, 0, \sin\ang{50}) =
(-0.643, 0, 0.766)$. Then
$\overrightarrow{AA'} = 400(0.866\,\mathbf s - 0.5\,\mathbf u) =
(128.6, 346.4, -153.2)$ in metres. Its \omterm{def:b1:field-lab-methods-1:trend-plunge}{trend} is $\arctan(128.6/346.4) =
\ang{20.4}$ and its \omterm{def:b1:field-lab-methods-1:trend-plunge}{plunge} $\arcsin(153.2/400) = \ang{22.5}$, which is also
$\arcsin(\sin\ang{30}\sin\ang{50})$. (c) An \omterm{def:b1:faults-folds-fabrics:fault-types}{oblique-slip fault}: the \omterm{def:b1:faults-folds-fabrics:fault-elements}{rake} lies
between $0$ and \ang{-90}, closer to $0$: a \omterm{def:b1:faults-folds-fabrics:fault-types}{sinistral} fault with a normal
component.
\end{solution}

\begin{solution}{exo:b1:faults-folds-fabrics:6}
(a) The two faults make \ang{60} with each other; $\sigma_1$ bisects the acute
angle, so it is horizontal along $050$ (and $230$); $\sigma_3$ is horizontal
along $140$; $\sigma_2$ is vertical, along the line in which the two vertical
faults meet. Each fault makes \ang{30} with $\sigma_1$, so $\ang{45} - \phi/2 =
\ang{30}$ and $\phi = \ang{30}$. (b) The two faults cut the ground into four
wedges; the wedges that contain $\sigma_1$, toward $050$ and toward $230$, move
toward each other, and those that contain $\sigma_3$ move apart, toward $140$
and $320$. Take the fault striking $020$, between the wedge toward $050$ (to its
east) and the wedge toward $320$ (to its west): relative to the western wedge,
the eastern one moves toward $230$ and away from $320$, both of which have a
component toward $200$ along the fault. An observer on the western wedge,
facing east, sees the eastern one move toward $200$, to the right: the fault
striking $020$ is \omterm{def:b1:faults-folds-fabrics:fault-types}{dextral}. The same argument, or the symmetry about
$\sigma_1$, makes the fault striking $080$ \omterm{def:b1:faults-folds-fabrics:fault-types}{sinistral}.
\end{solution}

\begin{solution}{exo:b1:faults-folds-fabrics:7}
The upward normals are
\begin{align*}
  \mathbf n_1 &= (\sin\ang{40}\sin\ang{120}, \sin\ang{40}\cos\ang{120}, \cos\ang{40}) = (0.557, -0.321, 0.766),\\
  \mathbf n_2 &= (\sin\ang{30}\sin\ang{260}, \sin\ang{30}\cos\ang{260}, \cos\ang{30}) = (-0.492, -0.087, 0.866).
\end{align*}
Their cross product is
$(-0.211, -0.859, -0.207)$, which points downward: the \omterm{def:b1:faults-folds-fabrics:fold-elements}{hinge line} \omterm{def:b1:field-lab-methods-1:trend-plunge}{plunges}
$\arcsin 0.207/0.909 = \ang{13.1}$ toward $\ang{180} + \arctan(0.211/0.859) =
\ang{193.8}$. $\mathbf n_1\cdot\mathbf n_2 = 0.417$, the angle between the
normals is \ang{65.3}, and $\iota = \ang{114.7}$: an \omterm{def:b1:faults-folds-fabrics:fold-classes}{open fold}.
\end{solution}

\begin{solution}{exo:b1:faults-folds-fabrics:8}
$\mathbf n_1 = (0.321, 0.383, 0.866)$ for the bedding and
$\mathbf n_2 = (0.337, -0.925, 0.174)$ for the \omterm{def:b1:faults-folds-fabrics:foliation}{cleavage}. Their cross product,
taken downward, gives a line that \omterm{def:b1:field-lab-methods-1:trend-plunge}{plunges} \ang{25.4} toward \ang{074.8}. By
\cref{prop:b1:faults-folds-fabrics:intersection}, this is the direction of the
\omterm{def:b1:faults-folds-fabrics:fold-elements}{hinge line}: the \omterm{def:b1:geological-maps:fold}{fold} \omterm{def:b1:field-lab-methods-1:trend-plunge}{plunges} about \ang{25} toward the east-north-east, which
any other \omterm{def:b1:geological-maps:geological-map}{outcrop} of either limb should confirm.
\end{solution}

\begin{solution}{exo:b1:faults-folds-fabrics:9}
(a) An observer south of the fault, facing north, sees the northern part of the
dyke displaced to the east, to the right: the fault is \omterm{def:b1:faults-folds-fabrics:fault-types}{dextral}. (b) The dyke is
a vertical plane perpendicular to the fault, so its offset along the \omterm{def:b1:geological-maps:attitude}{strike} is
the strike-slip component: \qty{350}{m}, which is the whole slip if the
\omterm{def:b1:faults-folds-fabrics:fault-rocks}{slickenlines} are horizontal. (c) The slip then makes \ang{20} with the \omterm{def:b1:geological-maps:attitude}{strike}
in the vertical fault plane: its horizontal component is still \qty{350}{m},
so the \omterm{def:b1:faults-folds-fabrics:fault-elements}{net slip} is $350/\cos\ang{20} = \qty{372.5}{m}$ and its vertical
component $350\tan\ang{20} = \qty{127.4}{m}$. The northern block moved east,
and the \omterm{def:b1:faults-folds-fabrics:fault-rocks}{slickenlines} descend toward the east, so it also went down, by
\qty{127.4}{m}. A vertical dyke cannot show that vertical movement, since
sliding it along itself leaves its trace unchanged.
\end{solution}

\begin{solution}{exo:b1:faults-folds-fabrics:10}
(a) The normal to the fault makes $\ang{90} - \theta$ with $\sigma_1$, so by
\cref{prop:b1:stress-strain-rheology:mohr},
$\sigma_n' = \sigma_1'\sin^2\theta + \sigma_3'\cos^2\theta$ and
$\tau = (\sigma_1' - \sigma_3')\sin\theta\cos\theta$. Slip starts when
$\tau = \mu\sigma_n'$; dividing by $\sigma_3'\cos^2\theta$ and writing
$R = \sigma_1'/\sigma_3'$, $(R - 1)\tan\theta = \mu(R\tan^2\theta + 1)$, so
$R\tan\theta(1 - \mu\tan\theta) = \tan\theta + \mu$, which is the formula.
(b) With $\sigma_1$ vertical, $\theta$ is \ang{90} minus the \omterm{def:b1:geological-maps:attitude}{dip}:
$\theta = \ang{30}$ gives $R = 3.12$; $\theta = \ang{45}$, $R = 4.00$;
$\theta = \ang{55}$, $R = 9.92$. (c) The denominator vanishes when
$\tan\theta = 1/\mu$, $\theta = \ang{59.0}$; beyond it, $R$ would be negative:
no ratio of positive effective stresses can make the fault slip, however large
$\sigma_1'$ becomes. \omterm{def:b1:faults-folds-fabrics:fault-types}{Normal faults} dipping less than \ang{31.0} cannot slip
under a vertical $\sigma_1$ with $\mu = 0.6$. (d) The best-oriented plane is at
$\ang{45} - \phi/2$ from $\sigma_1$. For $\mu = 0.6$, $\phi = \ang{31.0}$:
\omterm{def:b1:faults-folds-fabrics:fault-types}{normal faults} \omterm{def:b1:geological-maps:attitude}{dip} \ang{60.5}, thrusts \ang{29.5}. For $\mu = 0.85$,
$\phi = \ang{40.4}$: \ang{65.2} and \ang{24.8}.
\end{solution}

\begin{solution}{exo:b1:faults-folds-fabrics:11}
Looking north, east is to the right, and planes dipping west rise toward the
right. (a) At B the bedding \omterm{def:b1:geological-maps:attitude}{dips} the same way as the \omterm{def:b1:faults-folds-fabrics:foliation}{cleavage}, more steeply:
by \cref{prop:b1:faults-folds-fabrics:cleavage-bedding}(a), B is on an
overturned limb; A and C, where the bedding is less steep, are the right way
up. (b) The \omterm{def:b1:faults-folds-fabrics:foliation}{cleavage} rises at \ang{60}; at A and C the bedding rises at
\ang{40}, and the acute rotation from one to the other is clockwise, through
\ang{20}, moving the top of the \omterm{def:b1:faults-folds-fabrics:foliation}{cleavage} to the east; at B the bedding rises
at \ang{80} and the rotation is anticlockwise, through \ang{20}, moving it to
the west. (c) The sense changes between A and B and again between B and C:
there is a \omterm{def:b1:faults-folds-fabrics:fold-elements}{hinge} in each interval. At A, which is the right way up, the top of
the \omterm{def:b1:faults-folds-fabrics:foliation}{cleavage} moves east, toward the \omterm{def:b1:geological-maps:fold}{anticline} at the top of the limb: the \omterm{def:b1:faults-folds-fabrics:fold-elements}{hinge}
between A and B is an \omterm{def:b1:geological-maps:fold}{anticline}, and the one between B and C a \omterm{def:b1:geological-maps:fold}{syncline}.
(d) In the notation of the proof, the right-way-up \omterm{def:b1:faults-folds-fabrics:fold-elements}{limbs} \omterm{def:b1:geological-maps:attitude}{dip}
$\gamma - \iota/2 = \ang{40}$ and the overturned \omterm{def:b1:faults-folds-fabrics:fold-elements}{limbs} $\gamma + \iota/2 =
\ang{80}$, so $\gamma = \ang{60}$, as the \omterm{def:b1:faults-folds-fabrics:foliation}{cleavage} says, and $\iota =
\ang{40}$: \omterm{def:b1:faults-folds-fabrics:fold-classes}{tight folds}. The \omterm{def:b1:faults-folds-fabrics:fold-elements}{axial planes} \omterm{def:b1:geological-maps:attitude}{strike} north and \omterm{def:b1:geological-maps:attitude}{dip} \ang{60} to the
west, \omterm{def:b1:geological-maps:attitude}{attitude} $270/60$. The short \omterm{def:b1:faults-folds-fabrics:fold-elements}{limbs} have turned past the vertical with
their tops toward the east: the \omterm{def:b1:geological-maps:fold}{folds} verge to the east, as in the figure of
the chapter.
\end{solution}

\begin{solution}{exo:b1:faults-folds-fabrics:12}
(a) In a symmetric chevron \omterm{def:b1:geological-maps:fold}{fold} each limb makes the angle $\iota/2$ with the
\omterm{def:b1:faults-folds-fabrics:fold-elements}{axial plane}, which is perpendicular to the direction of shortening; a limb of
length $l$ therefore spans $l\sin(\iota/2)$ across the \omterm{def:b1:geological-maps:fold}{folds}, and the whole
layer $L_0\sin(\iota/2)$. (b) $\sin\ang{30} = 0.5$: a shortening of
\qty{50}{\percent}. (c) $\sin(\iota/2) = 0.7$ gives $\iota = \ang{88.9}$.
(d) $\iota = \ang{156}$ and $\sin\ang{78} = 0.978$: a shortening of
\qty{2.2}{\percent}, which is also $1 - \cos\ang{12}$.
\end{solution}

\begin{solution}{pb:b1:faults-folds-fabrics:1}
\textbf{1.} $\mathbf n = (\sin\delta\sin\alpha, \sin\delta\cos\alpha,
\cos\delta)$ gives
\[
  \mathbf n_1 = (0.20855, 0.03415, 0.97742), \qquad
  \mathbf n_2 = (-0.20855, 0.03415, 0.97742).
\]

\textbf{2.} $\mathbf n_1\cdot\mathbf n_2 = -0.04349 + 0.00117 + 0.95535 =
0.91302$, and the angle between the normals is \ang{24.1}.

\textbf{3.} By \cref{prop:b1:faults-folds-fabrics:hinge},
$\iota = \ang{180} - \ang{24.1} = \ang{155.9}$.

\textbf{4.} Between \ang{170} and \ang{90}: an \omterm{def:b1:faults-folds-fabrics:fold-classes}{open fold}.

\textbf{5.} The \omterm{def:b1:faults-folds-fabrics:fold-elements}{limbs} \omterm{def:b1:geological-maps:attitude}{dip} toward $080.7$ and $279.3$, not due east and west,
because the \omterm{def:b1:geological-maps:fold}{fold} \omterm{def:b1:field-lab-methods-1:trend-plunge}{plunges}: their steepest lines are not in the profile, which
shows only an \omterm{def:b1:cross-sections:apparent-dip}{apparent dip}. In an east--west section, by
\cref{thm:b1:cross-sections:apparent-dip}, $\tan\beta = \tan\ang{12.2}
\cos\ang{9.3} = 0.2134$, so $\beta = \ang{12.0}$, and $\ang{180} -
2\times\ang{12.0} = \ang{155.9}$ (the profile is tilted \ang{2} from that
section, which changes nothing to this precision).

\textbf{6.} Its \omterm{def:b1:faults-folds-fabrics:fold-elements}{limbs} are parallel to those of the \omterm{def:b1:geological-maps:fold}{anticline}, so its \omterm{def:b1:faults-folds-fabrics:fold-elements}{interlimb
angle} is also \ang{155.9}.

\textbf{7.} $\mathbf n_1\times\mathbf n_2 = (0.03415\times0.97742 -
0.97742\times0.03415,\; 0.97742\times(-0.20855) - 0.20855\times0.97742,\;
0.20855\times0.03415 + 0.03415\times0.20855) = (0, -0.40768, 0.01424)$.
Taken downward, $(0, 0.40768, -0.01424)$: the \omterm{def:b1:faults-folds-fabrics:fold-elements}{hinge line} \omterm{def:b1:field-lab-methods-1:trend-plunge}{trends} $000$ and
\omterm{def:b1:field-lab-methods-1:trend-plunge}{plunges} $\arctan(0.01424/0.40768) = \ang{2.0}$ to the north.

\textbf{8.} The \omterm{def:b1:faults-folds-fabrics:fold-elements}{axial plane} contains the \omterm{def:b1:faults-folds-fabrics:fold-elements}{hinge line} and bisects the two \omterm{def:b1:faults-folds-fabrics:fold-elements}{limbs};
its normal is along $\mathbf n_1 - \mathbf n_2 = (0.41710, 0, 0)$, horizontal
and east--west. It is vertical and \omterm{def:b1:geological-maps:attitude}{strikes} north--south: $090/90$. The
\omterm{def:b1:geological-maps:fold}{anticline} is upright.

\textbf{9.} An open, \omterm{def:b1:faults-folds-fabrics:fold-classes}{upright fold}, plunging \ang{2.0} to the north.

\textbf{10.} By \cref{prop:b1:geological-maps:folds}, the \omterm{def:b1:geological-maps:geological-map}{outcrops} of an
\omterm{def:b1:geological-maps:fold}{anticline} close in the direction of \omterm{def:b1:field-lab-methods-1:trend-plunge}{plunge}: to the north. Its core holds the
oldest rocks: the sandstone at the base of the cover and, where the river has
cut deepest, the basement of schist and granite beneath it.

\textbf{11.} Axial-planar, the \omterm{def:b1:faults-folds-fabrics:foliation}{cleavage} is vertical and \omterm{def:b1:geological-maps:attitude}{strikes} north--south,
$090/90$. By \cref{prop:b1:faults-folds-fabrics:intersection}, its
\omterm{def:b1:faults-folds-fabrics:lineation}{intersection lineation} with the bedding is parallel to the \omterm{def:b1:faults-folds-fabrics:fold-elements}{hinge line}: \omterm{def:b1:field-lab-methods-1:trend-plunge}{trend}
$000$, \omterm{def:b1:field-lab-methods-1:trend-plunge}{plunge} \ang{2.0}. Directly: the plane $080.7/12.2$ meets the vertical
plane striking north along a line of \omterm{def:b1:field-lab-methods-1:trend-plunge}{plunge} $\arctan(\tan\ang{12.2}
\cos\ang{80.7}) = \ang{2.0}$ toward $000$.

\textbf{12.} Looking north, the east limb rises toward the left at \ang{12.0}
and the \omterm{def:b1:faults-folds-fabrics:foliation}{cleavage} is vertical: the acute rotation from the \omterm{def:b1:faults-folds-fabrics:foliation}{cleavage} to the
bedding, through \ang{78.0}, is anticlockwise and moves the top of the
\omterm{def:b1:faults-folds-fabrics:foliation}{cleavage} to the west. The limb is the right way up, so by
\cref{prop:b1:faults-folds-fabrics:cleavage-bedding}(b) the \omterm{def:b1:geological-maps:fold}{anticline} \omterm{def:b1:faults-folds-fabrics:fold-elements}{hinge} is
to the west, at $x = \qty{3.5}{km}$, as the map shows.

\textbf{13.} The \omterm{def:b1:geological-maps:attitude}{strike} is $120 - 90 = 030$; the plane \omterm{def:b1:geological-maps:attitude}{dips} toward $120$, so
the \omterm{def:b1:faults-folds-fabrics:fault-elements}{hanging wall} lies on the east-south-east side of the trace.

\textbf{14.} The fault is normal: $\sigma_1$ was vertical.

\textbf{15.} $\sigma_2$ is horizontal along the \omterm{def:b1:geological-maps:attitude}{strike}, $030$--$210$;
$\sigma_3$ horizontal across it, $120$--$300$.

\textbf{16.} $\ang{45} + \phi/2 = \ang{60}$ gives $\phi = \ang{30}$, and
$\mu = \tan\ang{30} = 0.577$.

\textbf{17.} The conjugate set \omterm{def:b1:geological-maps:attitude}{strikes} $030$ as well and \omterm{def:b1:geological-maps:attitude}{dips} \ang{60} the other
way: $300/60$. The two sets meet at $\ang{180} - 2\times\ang{60} = \ang{60}$,
with the vertical $\sigma_1$ bisecting the acute angle.

\textbf{18.} With $\mu = 0.85$, $\phi = \ang{40.4}$ and the best-oriented
\omterm{def:b1:faults-folds-fabrics:fault-types}{normal faults} \omterm{def:b1:geological-maps:attitude}{dip} \ang{65.2}; the Varden fault, at \ang{60}, is \ang{5.2} off.
By \cref{exo:b1:faults-folds-fabrics:10}, it slips at the ratio
$\sigma_1'/\sigma_3' = (1 + 0.85\cot\ang{30})/(1 - 0.85\tan\ang{30}) = 4.85$,
against $4.68$ for the best plane: it is nearly, but not exactly, the best
oriented, and would slip again at a stress ratio only \qty{4}{\percent}
higher.

\textbf{19.} The slip is down the \omterm{def:b1:geological-maps:attitude}{dip}: $r = \ang{-90}$; the \omterm{def:b1:faults-folds-fabrics:fault-elements}{slip vector} \omterm{def:b1:field-lab-methods-1:trend-plunge}{trends}
$120$ and \omterm{def:b1:field-lab-methods-1:trend-plunge}{plunges} \ang{60}.

\textbf{20.} $(\cos\ang{60}\sin\ang{120}, \cos\ang{60}\cos\ang{120},
-\sin\ang{60}) = (0.4330, -0.2500, -0.8660)$.

\textbf{21.} $e = 200/\tan\ang{60} = \qty{115.5}{m}$.

\textbf{22.} East: $115.47\sin\ang{120} = \qty{100.0}{m}$; north:
$115.47\cos\ang{120} = \qty{-57.7}{m}$, that is \qty{57.7}{m} to the south.

\textbf{23.} $L = 200/\sin\ang{60} = \qty{230.9}{m}$, and
$200^2 + 115.47^2 = 40\,000 + 13\,333 = 53\,333 = 230.94^2$. The \omterm{def:b1:cross-sections:throw}{throw} and the
\omterm{def:b1:cross-sections:throw}{heave} are the vertical and horizontal components of the slip, which lies in
the vertical plane of the \omterm{def:b1:geological-maps:attitude}{dip}: they are the sides of a right triangle whose
hypotenuse is the slip.

\textbf{24.} To the nearest metre, the \omterm{def:b1:faults-folds-fabrics:fault-elements}{net slip} on the Varden fault is
\textbf{\qty[text-series-to-math=true]{231}{m}}.
\end{solution}
